% Backend: mineru (hybrid-auto-engine)
% Raw content: raw/input/hybrid_auto/input.md
% Source: Shao, Mathematical Statistics, §3.2.2 (page 192)
% agent_corrected: true (env_wrap from plain-text "Theorem 3.4" header)
\documentclass{article}
\usepackage{amsmath,amssymb,amsthm}
\newtheorem{theorem}{Theorem}
\begin{document}

\section*{3.2.2 Variances of U-statistics}

If $E[h(X_1,\dots,X_m)]^2 < \infty$, then the variance of $U_n$ in (3.11) with kernel $h$
has an explicit form. To derive $\mathrm{Var}(U_n)$, we need some notation. For $k=1,\dots,m$, let
\[
  h_k(x_1,\dots,x_k) = E[h(X_1,\dots,X_m) \mid X_1=x_1,\dots,X_k=x_k]
                     = E[h(x_1,\dots,x_k,X_{k+1},\dots,X_m)].
\]
Note that $h_m = h$. It can be shown that
\[
  h_k(x_1,\dots,x_k) = E[h_{k+1}(x_1,\dots,x_k,X_{k+1})]. \tag{3.14}
\]
Define
\[
  \tilde{h}_k = h_k - E[h(X_1,\dots,X_m)], \tag{3.15}
\]
$k=1,\dots,m$, and $\tilde{h}=\tilde{h}_m$. Then, for any $U_n$ defined by (3.11),
\[
  U_n - E(U_n) = \binom{n}{m}^{-1} \sum_{c} \tilde{h}(X_{i_1},\dots,X_{i_m}). \tag{3.16}
\]

\begin{theorem}[Hoeffding's theorem]
\label{thm:hoeffding-u-statistic}
For a U-statistic $U_n$ given by (3.11) with $E[h(X_1,\dots,X_m)]^2 < \infty$,
\[
  \mathrm{Var}(U_n) = \binom{n}{m}^{-1} \sum_{k=1}^{m} \binom{m}{k} \binom{n-m}{m-k} \zeta_k,
\]
where
\[
  \zeta_k = \mathrm{Var}(h_k(X_1,\dots,X_k)).
\]
\end{theorem}

\begin{proof}
Consider two sets $\{i_1,\dots,i_m\}$ and $\{j_1,\dots,j_m\}$ of $m$ distinct
integers from $\{1,\dots,n\}$ with exactly $k$ integers in common. The number of
distinct choices of two such sets is $\binom{n}{m}\binom{m}{k}\binom{n-m}{m-k}$.
By the symmetry of $\tilde{h}_m$ and independence of $X_1,\dots,X_n$,
\[
  E[\tilde{h}(X_{i_1},\dots,X_{i_m}) \tilde{h}(X_{j_1},\dots,X_{j_m})] = \zeta_k \tag{3.17}
\]
for $k=1,\dots,m$ (exercise). Then, by (3.16),
\[
  \mathrm{Var}(U_n)
    = \binom{n}{m}^{-2} \sum_c \sum_c E[\tilde{h}(X_{i_1},\dots,X_{i_m}) \tilde{h}(X_{j_1},\dots,X_{j_m})]
    = \binom{n}{m}^{-2} \sum_{k=1}^{m} \binom{n}{m} \binom{m}{k} \binom{n-m}{m-k} \zeta_k.
\]
This proves the result.
\end{proof}

\end{document}
